\section{Discussion}\label{sec:discussion}

\subsection{Coverage is not test quality}
Branch coverage and mutation score diverge in our own data at the sample level --- exactly
why we report both for each suite in RQ1/RQ2. On \texttt{JA-008}, the baseline suite
reaches 60\% branch coverage yet produces a 0\% mutation score: it exercises the
function's branches but makes no assertion about any returned value, correct or not. It is
a poorly performing suite, and any reading of RQ1/RQ2 must acknowledge as much: the
effective-subset coverage figures (75.0\% Python, 50.0\% Java) sit far above their
mutation-score counterparts (36.84\% Python, 0\% Java). Given coverage alone, one might
naively call gpt-4o-mini's tests somewhat effective; given mutation scores alone, one
might call them largely ineffective. Both readings are correct in aggregate, which is why
we report both figures at every level.

\subsection{Where does the ``breaking point'' actually lie?}
GAP-T asks for the cyclomatic-complexity threshold beyond which LLM test quality begins to
degrade, but RQ3 finds no such threshold ($\rho \in [+0.13, +0.18]$, $p>0.4$ for either
language) anywhere inside the 5--10 CC band. Read next to RQ1/RQ2, what dominates instead
is a binary outcome: whether the suite compiles and, if it does, whether it exercises the
sampled method rather than another with the same name or an overloaded parent method. A
CC~5 function and a CC~10 function in our sample are equally likely to fail that gate.
Complexity may well drive test-generation quality outside this range, but inside the
narrow band we sampled --- chosen to exclude the extremes --- something else dominates the
variance. Identifying a true ``breaking point'' would require sampling much lower- and
higher-CC functions as controls, which we defer to Section~\ref{sec:conclusion}.

\subsection{Cost is not the primary constraint; correctness is}
All 120 one-shot suites fit within a \$1 budget (the pre-registered estimate;
amendment~v1.1, \S8.2) and appeared within minutes, whereas a 60-second search budget per
class left EvoSuite and Randoop running for tens of minutes before exercising any of the
60 Java functions. On both price and speed, the LLM (\texttt{gpt-4o-mini-2024-07-18})
sits far above the Java baselines. On the actual ability to exercise and kill mutants on
the intended target, however, it sits below both, with an effect size of $r=-1.00$: every
non-tied pair favours the automated Java tools. Cheap and fast but unreliable, the LLM's
true advantage lies in the time it saves --- and the interventions that time makes
possible. If part of that saved time were funnelled into suite reconfiguration (a
feedback loop, see below), it might well close the gap without approaching the baselines'
total runtime.

\subsection{The \textsc{invalid} rate as a finding, not a bug}
We made no effort to exclude functions for which a one-shot external test is structurally
hard --- private target methods above all, of which there were 17/60 Java in our sample.
Filtering such functions after discovering Java's low effective rate would introduce a
selection effect and violate the pre-registered protocol (amendments~v1.1/v1.2,
Section~\ref{sec:threats}). And the data speak plainly to another cause: the primary
reason for Java's low test quality was wrong-target invocation. A private target is
involved in only 1 of the 9 Java functions our pipeline scored \textsc{invalid} and only 14
of the 46 that compiled but touched nothing (Section~\ref{sec:results}). The outright compile
failures are distributed differently, concentrated in the API-heavy \texttt{AccurateMath}
and \texttt{Gson} classes (Section~\ref{sec:results}).

This looks like a general limitation of one-shot, zero-context test generation for Java,
not something specific to our sample. Haratian et al.~\cite{haratian2026praware} report a
63\% compilation-error rate for GPT-4o across 353 real Java pull requests --- roughly
two-thirds of our Java \textsc{invalid}-or-no-touch rate (91.7\%), and another one-shot,
zero-context task on real, uncurated code. Sapozhnikov
et al.~\cite{sapozhnikov2024empirical} report over 50\% malformed or non-compiling
ChatGPT-generated Java tests, and Kumar et al.~\cite{kumar2025empirical} tie weak
single-pass results to the absence of a repair loop. The interventions their literature
converges on --- feeding the model compiler or test-failure output
(Zeng et al.~\cite{zeng2026logic}; Xue et al.~\cite{xue2025distinct}), retrieving real API
definitions that target exactly the wrong-target problem we
observe~\cite{liu2025typeaware}, or sampling several candidate suites and keeping whichever
compiles --- each step outside the one-shot setting we evaluate. Our numbers are a lower
bound on what gpt-4o-mini can do with this codebase, and a direct demonstration of why the
field is moving towards feedback- and retrieval-augmented generation.

\subsection{Why gpt-4o-mini ``beats'' Pynguin}
The RQ2-B result against Pynguin ($r=+0.85$, in gpt-4o-mini's favour) speaks more to the
baseline than to the treatment: Pynguin's own median Python coverage on this dataset is
0.0\% (Section~\ref{sec:threats}). This stems from Pynguin's internal timeouts and a
\texttt{dill} serialisation failure on modules that import \texttt{\_json} --- neither of
which touches gpt-4o-mini's ability to write tests. We report the comparison anyway,
because a Python baseline was pre-registered (amendment~v1.2) and quietly dropping an
inconvenient result would itself be a validity threat. It is not a claim about
gpt-4o-mini's Python ability --- RQ1/RQ2-A already speak to that.

\subsection{Comparison with prior work}
Our coverage figures sit far below the headline numbers in this literature, and the gap is
informative rather than embarrassing. Wang et al.\ report up to 98.65\% line and 97.16\%
branch coverage for GPT-4o on TestEval~\cite{wang2025testeval}, and Kumar et al.\ report
99.05\% branch coverage for Gemini~\cite{kumar2025empirical}; our all-level median is 0\% and
our best effective median is 75.0\%. The decisive difference is the target, not the model:
those studies evaluate self-contained LeetCode-style functions, whereas every function here
is embedded in a real project and must be imported, constructed, and invoked in place. The
one prior result measured on real projects rather than puzzles agrees with us --- Haratian et
al.\ report GPT-4o passing only 37\% of Java pull requests against 36\% for
EvoSuite~\cite{haratian2026praware}, the same order of magnitude as our own effective rates.

Our mutation figures, by contrast, line up closely with prior work. We measure a 36.84\%
effective median for Python, against the 33.82\% aggregate reported by Sapozhnikov et
al.~\cite{sapozhnikov2024empirical} and the 33.8\% of TestForge~\cite{jain2025testgeneval}.
That agreement is worth stating explicitly: where coverage numbers diverge by 60 percentage
points across benchmarks, mutation numbers converge near one third. This is consistent with
the view that coverage is highly sensitive to how the target is framed while fault detection
is not, and it is the empirical core of our argument in Section~\ref{sec:results}.

Finally, our Java compile-failure rate is in the same band as the literature. Zhang et al.\
measure a 47.9--55.9\% Java compile-error rate on TestBench~\cite{zhang2024testbench}; our
direct \texttt{javac} audit puts 43/60 (71.7\%) of Java suites as failing to compile. Ours is
higher, which we attribute to sampling functions rather than whole classes: a class-level
prompt supplies the surrounding API for free, and a function-level prompt does not.

\subsection{Implications for practice}
Three practical conclusions follow, and they are narrower than ``LLMs can or cannot write
tests''.

\textbf{Do not substitute a one-shot LLM for a search-based generator on Java.} On matched
functions EvoSuite and Randoop killed strictly more mutants ($p=0.00044$ and $p=0.0033$, both
surviving Bonferroni correction). A team already running EvoSuite gains nothing by replacing
it with one-shot prompting at this complexity level.

\textbf{Treat the generated suite as a draft that must be executed before it is trusted.}
The dominant failure is not a weak assertion but a suite that never reaches the target: 46.7\%
of Python suites were invalid and 91.7\% of Java suites were invalid or touched nothing. A
developer who reads such a suite without running it will see plausible-looking test code that
verifies nothing. Any adoption should therefore gate on execution and, ideally, on a
fault-based check rather than on coverage.

\textbf{Give the model the surrounding API, not just the function.} Our exploratory RQ4/RQ5
results (below) and the contrast with class-level studies both point the same way: what the
model lacks is the context needed to construct the receiver, not the ability to write
assertions. For a team, that means the cheapest improvement is better prompt context, not a
larger model.

\subsection{RQ4 (exploratory, post-hoc): does real API context help?}
\label{sec:rq4}
Wrong-target invocation, the leading cause of Java invalidity above, motivated a small
post-hoc, exploratory intervention: the same 120 functions, the same
model/temperature/exemplar, but with the target class's (Java) or module's (Python) real
public-API skeleton injected from the pinned source with \texttt{javalang}/\texttt{ast}.
The new suites were generated in a separate output directory with their own results file,
so as not to interfere with the confirmatory $N=120$ analysis. The idea came only after we
had seen the $N=120$ results, so it sits outside the pre-registered RQ1--RQ3 and is
reported separately, to avoid HARKing.

\textbf{Java.} The number of suites our pipeline scored as compiling stays the same (51/60
in both arms; see Section~\ref{sec:threats} for why that flag overstates the true compile
rate), but the effective rate --- suites that touch the target --- more than doubles:
11/60 (18.3\%) with context versus 5/60 (8.3\%) without. Paired by function, 7/60 improve, 1/60 regresses, and 52/60
stay the same; a one-sided Wilcoxon test barely misses significance at $\alpha=0.05$, but
7 favourable pairs against 1 unfavourable is a noticeable trend, and the effect size among
the non-tied pairs is moderate-to-large. Mutation score does not improve at all ($n=48$
pairs; both medians 0\%). Context repairs the structural gap but not the behavioural one:
the model still fails to assert the expected return value, only that the method can be
invoked. The two problems need different interventions --- wrong-target invocation is an
API-navigation problem, weak assertions a test-design problem. Injecting a class skeleton
addresses the former; examples of valid assertions for the function at hand would address
the latter.

\textbf{Python.} Context appears to hurt under the pre-registered, file-level
all-or-nothing pass criterion: 7/60 files compiled with context versus 32/60 without. This
is almost certainly a measurement artefact of that criterion interacting with the number
of tests per file, not an actual drop in quality. Counted per test rather than per file,
the two arms are nearly identical: 267 test cases at 33.3\% passing with context, versus
233 test cases at 32.2\% without. What context actually does in Python is push gpt-4o-mini
to write more test cases per file --- most likely because it now sees more of the API
surface --- and the per-file, all-tests-must-pass criterion penalises those extra tests by
failing the whole file if any one of them trips over an unrelated behavioural detail. We
traced one such case, \texttt{flask.logging.has\_level\_handler}: an identical
misunderstanding of logger-hierarchy propagation appears in both arms, unrelated to the
presence or absence of context. This is a second-order validity threat worth flagging for
future work: an all-or-nothing file-level criterion penalises suite thoroughness by
turning one behavioural assertion error into an automatic file-level failure, so arms with
nearly identical per-test correctness can look dramatically different depending on how many
extra tests the model wrote. It is unlikely to affect future iterations that move to
per-test metrics, but it is worth noting as a source of noise.

\textbf{Implication.} Static API context is mostly beneficial: it improves Java's
wrong-target problem with no measurable downside, and Python's per-test performance is
essentially unaffected once measured fairly. Neither result changes the RQ1--RQ3
conclusions (Section~\ref{sec:results}): with or without context, gpt-4o-mini's one-shot
tests fall well short of the 80\%/60\% thresholds. Context costs little (no model calls
beyond the initial one-shot budget) and, for Java, is a real improvement worth adding to a
test-generation pipeline --- but including or excluding it does not change the overall
conclusion that one-shot generation is inadequate to the task.

\subsection{RQ5 (exploratory, post-hoc): one feedback/repair iteration on top of RQ4}
\label{sec:rq5}
We committed to a single repair iteration before attempting it, and iterated exactly once:
repeating a repair technique until some threshold is crossed is precisely the adaptation
our protocol forbids. Every RQ4 suite that was not yet effective was sent back to
gpt-4o-mini with a real failure signal and an explicit request to fix one failing aspect
of the suite. Java received the 0\%-coverage signal plus the target API skeleton
(Maven/JVM error messages proved too noisy to use as diagnostics); Python received the
fresh \texttt{pytest} failure output. Suites already effective at RQ4 (11/60 Java, and
Python's per-test baseline) were carried forward unchanged rather than regenerated.

\textbf{Java: further, if inconclusive, improvement.} The effective rate rises again, from
11/60 (18.3\%, RQ4) to 15/60 (25.0\%). Among the functions that changed, 4 improved and
none regressed: repair is at least as good as context across the board. The one-sided
Wilcoxon test barely misses significance again at $p=0.063$, but with an improved
rank-biserial effect size of $r=1.00$ (every non-tied pair favours repair) and a steady
climb across all three stages (8.3\% $\to$ 18.3\% $\to$ 25.0\%). Mutation score stays at a
median of 0\% ($n=45$ paired functions), the same as at RQ4. Repair helps gpt-4o-mini
reach the correct lines more often; just as with context, the assertions are still not
strong enough to kill mutants.

\textbf{Python: an actual decrease in performance, this time with no artefact to blame.}
The per-test-case pass rate falls to 15.5\% (106/682 tests), below both RQ4's 33.3\% and
the first context-less run's 32.2\%. File count and suite length are no longer confounders,
as the RQ4-vs-original test-case numbers show; unlike the context-vs-no-context comparison,
they do not explain this drop. Suite size does grow again --- to 682 tests across 60 files,
up from 267 at RQ4 --- but gpt-4o-mini appears to react to a visible failure signal by
expanding the suite rather than pinpointing the one faulty assertion and repairing it. The
additional tests fail at a higher rate than the initial ones succeeded, but this pattern
has no obvious mechanism or cause. We have no reason to believe it is a universal property
of gpt-4o-mini, of single-shot repair, or of the failure signal itself, but the
one-attempt limitation rules out further investigation.

\textbf{Overall assessment of RQ4+RQ5.} Two low-cost, non-fine-tuning interventions leave
Java with a real, if inconclusive, increase in coverage effectiveness and no downside
anywhere, and Python with no benefit from context plus an outright decrease from
single-shot repair. Neither intervention closes the gap to the 80\%/60\% thresholds, and
their uneven effects across the two languages suggest that ``add more scaffolding'' is not
a reliable way to improve LLM test generation in general. We pre-committed to stop
searching for such methods once two of them failed to meet the thresholds, and we do so
here rather than attempt a third.
